%% =================================================================
%% Kingma, Ba. "Adam: A Method for Stochastic Optimization."
%% ICLR 2015. arXiv:1412.6980v9.
%%
%% Reconstruction of page 2 (printed p. 2) as study notes.
%% Generated by: reproduce-all-pages-basic-tex (batch 1-15)
%%
%% Structure: Algorithm 1 box (Adam pseudocode) + body text
%% explaining the algorithm + Section 2.1 "Adam's update rule" (cut
%% off mid-sentence at the bottom). No numbered equations — all
%% math is inline or inside the algorithm box.
%% =================================================================

\documentclass[11pt]{article}

\usepackage[margin=1in]{geometry}
\usepackage{amsmath, amssymb}
\usepackage{mathrsfs}
\usepackage{bm}
\usepackage{xcolor}
\usepackage{fancyhdr}

\pagestyle{fancy}
\fancyhf{}
\lhead{\small Published as a conference paper at ICLR 2015}
\rhead{}
\cfoot{\small 2 $\to$ \arabic{page}}
\renewcommand{\headrulewidth}{0.4pt}

\setcounter{page}{1}

\begin{document}

% ---------- Algorithm 1 box ----------
\noindent
\fbox{\parbox{0.95\textwidth}{\small
\textbf{Algorithm 1:} \emph{Adam}, our proposed algorithm for
stochastic optimization. See section~2 for details, and for a
slightly more efficient (but less clear) order of computation.
$g_t^2$ indicates the elementwise square $g_t \odot g_t$. Good
default settings for the tested machine learning problems are
$\alpha = 0.001$, $\beta_1 = 0.9$, $\beta_2 = 0.999$ and
$\epsilon = 10^{-8}$. All operations on vectors are element-wise.
With $\beta_1^t$ and $\beta_2^t$ we denote $\beta_1$ and $\beta_2$
to the power $t$.\\[0.6em]
\textbf{Require:} $\alpha$: Stepsize\\
\textbf{Require:} $\beta_1, \beta_2 \in [0, 1)$: Exponential decay rates for the moment estimates\\
\textbf{Require:} $f(\theta)$: Stochastic objective function with parameters $\theta$\\
\textbf{Require:} $\theta_0$: Initial parameter vector\\[0.2em]
\hspace*{1em}$m_0 \leftarrow 0$ \hfill (Initialize $1^{\text{st}}$ moment vector)\\
\hspace*{1em}$v_0 \leftarrow 0$ \hfill (Initialize $2^{\text{nd}}$ moment vector)\\
\hspace*{1em}$t \leftarrow 0$ \hfill (Initialize timestep)\\
\hspace*{1em}\textbf{while} $\theta_t$ not converged \textbf{do}\\
\hspace*{2em}$t \leftarrow t + 1$\\
\hspace*{2em}$g_t \leftarrow \nabla_\theta f_t(\theta_{t-1})$ \hfill (Get gradients w.r.t.\ stochastic objective at timestep $t$)\\
\hspace*{2em}$m_t \leftarrow \beta_1 \cdot m_{t-1} + (1 - \beta_1) \cdot g_t$ \hfill (Update biased first moment estimate)\\
\hspace*{2em}$v_t \leftarrow \beta_2 \cdot v_{t-1} + (1 - \beta_2) \cdot g_t^2$ \hfill (Update biased second raw moment estimate)\\
\hspace*{2em}$\widehat{m}_t \leftarrow m_t / (1 - \beta_1^t)$ \hfill (Compute bias-corrected first moment estimate)\\
\hspace*{2em}$\widehat{v}_t \leftarrow v_t / (1 - \beta_2^t)$ \hfill (Compute bias-corrected second raw moment estimate)\\
\hspace*{2em}$\theta_t \leftarrow \theta_{t-1} - \alpha \cdot \widehat{m}_t / (\sqrt{\widehat{v}_t} + \epsilon)$ \hfill (Update parameters)\\
\hspace*{1em}\textbf{end while}\\
\hspace*{1em}\textbf{return} $\theta_t$ (Resulting parameters)
}}

\vspace{1em}

% ---------- Body prose ----------
\noindent
In section~2 we describe the algorithm and the properties of its
update rule. Section~3 explains our initialization bias correction
technique, and section~4 provides a theoretical analysis of Adam's
convergence in online convex programming. Empirically, our method
consistently outperforms other methods for a variety of models and
datasets, as shown in section~6. Overall, we show that Adam is a
versatile algorithm that scales to large-scale high-dimensional
machine learning problems.

\section*{2 \textsc{Algorithm}}

See algorithm~1 for pseudo-code of our proposed algorithm
\emph{Adam}. Let $f(\theta)$ be a noisy objective function: a
stochastic scalar function that is differentiable w.r.t. parameters
$\theta$. We are interested in minimizing the expected value of
this function, $\mathbb{E}[f(\theta)]$ w.r.t. its parameters
$\theta$. With $f_1(\theta), \ldots, f_T(\theta)$ we denote the
realisations of the stochastic function at subsequent timesteps
$1, \ldots, T$. The stochasticity might come from the evaluation
at random subsamples (minibatches) of datapoints, or arise from
inherent function noise. With
$g_t = \nabla_\theta f_t(\theta)$ we denote the gradient, i.e.
the vector of partial derivatives of $f_t$, w.r.t $\theta$
evaluated at timestep $t$.

The algorithm updates exponential moving averages of the gradient
($m_t$) and the squared gradient ($v_t$) where the hyper-parameters
$\beta_1, \beta_2 \in [0, 1)$ control the exponential decay rates of
these moving averages. The moving averages themselves are estimates
of the $1^\text{st}$ moment (the mean) and the $2^\text{nd}$ raw
moment (the uncentered variance) of the gradient. However, these
moving averages are initialized as (vectors of) 0's, leading to
moment estimates that are biased towards zero, especially during
the initial timesteps, and especially when the decay rates are
small (i.e.\ the $\beta$s are close to 1). The good news is that
this initialization bias can be easily counteracted, resulting in
bias-corrected estimates $\widehat{m}_t$ and $\widehat{v}_t$. See
section~3 for more details.

Note that the efficiency of algorithm~1 can, at the expense of
clarity, be improved upon by changing the order of computation,
e.g.\ by replacing the last three lines in the loop with the
following lines:
$\alpha_t = \alpha \cdot \sqrt{1 - \beta_2^t}/(1 - \beta_1^t)$
and
$\theta_t \leftarrow \theta_{t-1} - \alpha_t \cdot m_t/(\sqrt{v_t} + \hat{\epsilon})$.

\subsection*{2.1 \textsc{Adam's update rule}}

An important property of Adam's update rule is its careful choice
of stepsizes. Assuming $\epsilon = 0$, the effective step taken in
parameter space at timestep $t$ is
$\Delta_t = \alpha \cdot \widehat{m}_t / \sqrt{\widehat{v}_t}$.
The effective stepsize has two upper bounds:
$|\Delta_t| \leq \alpha \cdot (1 - \beta_1) / \sqrt{1 - \beta_2}$ in
the case $(1 - \beta_1) > \sqrt{1 - \beta_2}$, and
$|\Delta_t| \leq \alpha$

% [page ends here — continues on page 3]

\end{document}
